% La différence entre \arabic{compteur} et \thecompteur se remarque si compteur est « lié » à un autre compteur. Un compteur est lié à un autre si ce premier est réinitialisé lorsque le second est incrémenté. Par exemple dans la classe book, les figures sont numérotées en commençant par le numéro du chapitre. Ainsi, la deuxième figure du chapitre 3 est numérotée « 3.2 », et la valeur de \arabic{figure} est alors « 2 », tandis que celle de \thefigure est « 3.2 ».

%La différence étant que refstepcounter donne à la commande \ref{label} la même que \thecompteur (au lieu de \arabic{compteur}, vous suivez ?).

% informations provenant de : http://blog.dorian-depriester.fr/latex/les-compteurs-sous-latex#Personnaliser_lrsquoaffichage_du_compteur

\tcbuselibrary{theorems, skins, breakable} % Pour utiliser certaines commandes du package "tcolorbox"

\usetikzlibrary{shadows}


%%%%%%%%%%%%%%%%%%%%%%%%%%%%%%%%%%%%%%%%%%%%%%%%%%%%%%%%%%%%%%%%%%%%%%%%%%
%%%%%%%%%%%%%%%%%%%%%%%%%%%   ENVIRONNEMENTS   %%%%%%%%%%%%%%%%%%%%%%%%%%%
%%%%%%%%%%%%%%%%%%%%%%%%%%%%%%%%%%%%%%%%%%%%%%%%%%%%%%%%%%%%%%%%%%%%%%%%%%
%-------------------
%---- Théorèmes ----
%-------------------
\newcounter{compteurthm}[chapter] % Ici on définit un nouveau compteur associé au compteur "chapter", i.e. quand "chapter" augmente de niveau, "compteurthm" est remis à zéro.
\renewcommand{\thecompteurthm}{\thechapter.\arabic{compteurthm}} % On redéfini \thecompteurthm pour que sa valeur soit précédée du numéro de chapitre.
\newenvironment{thm}[1]
	{
		\refstepcounter{compteurthm}         
         \begin{center}%
            \begin{tikzpicture}%
                \node[rectangle, draw=black, rounded corners=0pt, inner xsep=5pt, inner ysep=6pt, outer ysep=10pt]\bgroup                     
                		\begin{minipage}{#1\linewidth}
                			{\bf Théorème \thecompteurthm\hspace{2pt}:}
                			
                			\vspace{2mm}\hspace{7.5mm}
                			\begin{minipage}{0.93\linewidth}\it
                			

    }
	{                		
                			\end{minipage}
                		\end{minipage}\egroup;%
            \end{tikzpicture}%
         \end{center}%
	}




%--------------------------
%---- Théorèmes nommés ----
%--------------------------
\newenvironment{thmnom}[2]
	{
		\refstepcounter{compteurthm}         
         \begin{center}%
            \begin{tikzpicture}%
                \node[rectangle, draw=black, rounded corners=0pt, inner xsep=5pt, inner ysep=6pt, outer ysep=10pt]\bgroup                     
                		\begin{minipage}{#1\linewidth}
                			{\bf Théorème \thecompteurthm\hspace{2pt} : \textsl{(#2)}}
                			
                			\vspace{2mm}\hspace{7.5mm}
                			\begin{minipage}{0.93\linewidth}\it

    }
	{                		
                			\end{minipage}
                		\end{minipage}\egroup;%
            \end{tikzpicture}%
         \end{center}%
	}











%--------------------------------------
%---- Personalisation de tcolorbox ----
%--------------------------------------
\tcbset{
	semiboxedshadow/.style 2 args={
			enhanced, colback=gray!20, colframe=gray!10,rounded corners, arc=3mm, left=4mm,top=4mm, bottom=3mm,%boxrule=0pt,
			%borderline west={#1}{0pt}{#2},left={#1*14/10+7pt},top=1ex, bottom=1ex,
			%borderline south={#1}{0pt}{#2},left={#1*14/10+7pt},top=1ex, bottom=1ex,
			fuzzy shadow={-1.5mm}{-1.5mm}{0mm}{0.3mm}{fill=black,opacity=0.25},			
			%drop fuzzy shadow,
			after skip=5mm
			},
	leftlined/.style 2 args={
			blanker, breakable,
			borderline west={#1}{0pt}{#2},left={#1*14/10+7pt},top=1ex, bottom=1ex,
			},
    exstyle/.style={enhanced, colback=Indigo!5, colframe=Indigo, 
           fonttitle=\bfseries, 
           colbacktitle=Indigo, coltitle=white,  
           top=2mm,
           boxed title style={},
           attach boxed title to top left={yshift=-\tcboxedtitleheight/2,xshift=4mm}%
           },
    teostyle/.style={enhanced, colback=white, colframe=White, 
           fonttitle=\bfseries, sharp corners, boxrule=0pt,
           colbacktitle=Red!100, coltitle=white,
           drop fuzzy shadow,  
           top=2mm,
           boxed title style={sharp corners},
           attach boxed title to top left={}%
           },
	indented/.style={
			breakable,left=#1, colback=white, colframe=white,
			%borderline west={#1}{0pt}{#2},left={#1*14/10+7pt},top=1ex, bottom=1ex,
			}, 
}


%\newtcbtheorem[number within=section]{miTeorema}{Teorema}{teostyle}{teo}

% CI DESSOUS UN EXEMPLE D'UTILISATION DU STYLE "THEOSTYLE" CI-DESSUS :
%\begin{miTeorema}{Un teorema}{label}
%Sean $a, b\in \mathbb{Z}$ con $b\neq 0$ . Sea $q\in \mathbb{Z}$ \dots
%\end{miTeorema}



%----------------------
%---- Définition-s ----
%----------------------
\newtheoremstyle{definition}% name
{3pt}% Space above  % -> NE MARCHE PAS - CERTAINEMENT A CAUSE DE TCB QUI DOIT "OVERRULER" LES DIMENSIONS
{3pt}% Space below  % -> NE MARCHE PAS - CERTAINEMENT A CAUSE DE TCB QUI DOIT "OVERRULER" LES DIMENSIONS
{}% Body font \itshape
{}% Indent amount
{\bfseries}% Theorem head font
{\vspace{1mm}\newline}% Punctuation after theorem head
{1mm}%{0pt\Needspace{3\baselineskip}}% Space after theorem head % -> NE MARCHE PAS - CERTAINEMENT A CAUSE DE TCB QUI DOIT "OVERRULER" LES DIMENSIONS
{}% Theorem head spec (can be left empty, meaning ‘normal’ )

\newtheoremstyle{definitions}% name
{3pt}% Space above  % -> NE MARCHE PAS - CERTAINEMENT A CAUSE DE TCB QUI DOIT "OVERRULER" LES DIMENSIONS
{3pt}% Space below  % -> NE MARCHE PAS - CERTAINEMENT A CAUSE DE TCB QUI DOIT "OVERRULER" LES DIMENSIONS
{}% Body font \itshape
{}% Indent amount
{\bfseries}% Theorem head font
{\vspace{1mm}\newline}% Punctuation after theorem head
{1mm}%{0pt\Needspace{3\baselineskip}}% Space after theorem head % -> NE MARCHE PAS - CERTAINEMENT A CAUSE DE TCB QUI DOIT "OVERRULER" LES DIMENSIONS
{}% Theorem head spec (can be left empty, meaning ‘normal’ )

%\tcolorboxenvironment{definition}{exstyle}
%\tcolorboxenvironment{definition}{teostyle}
\tcolorboxenvironment{definition}{semiboxedshadow}
\tcolorboxenvironment{definitions}{semiboxedshadow}


\theoremstyle{definition}
\newtheorem{definition}{Définition}[chapter]
\theoremstyle{definitions}
\newtheorem{definitions}[definition]{Définitions}






%----------------------
%---- Propositions ----
%----------------------
\newtheoremstyle{proposition}% name
{3pt}% Space above
{3pt}% Space below
{\itshape}% Body font
{}% Indent amount
{\bfseries}% Theorem head font
{\vspace{1mm}\newline}% Punctuation after theorem head
{1cm}%{0pt\Needspace{3\baselineskip}}% Space after theorem head
{}% Theorem head spec (can be left empty, meaning ‘normal’ )

\tcolorboxenvironment{proposition}{leftlined={5pt}{red!30!black}}
\theoremstyle{proposition}
\newtheorem{proposition}{Proposition}[chapter]
%\newtcbtheorem[number within=chapter]{proposition}{Proposition}
%\declaretheorem[style=proposition]{named}





%-------------------
%---- Théorèmes ----
%-------------------
\newtheoremstyle{theoreme}% name
{3pt}% Space above
{1cm}% Space below
{\itshape}% Body font
{}% Indent amount
{\bfseries}% Theorem head font
{\vspace{1mm}\newline}% Punctuation after theorem head
{1cm}%{0pt\Needspace{3\baselineskip}}% Space after theorem head
{}% Theorem head spec (can be left empty, meaning ‘normal’ )

\tcolorboxenvironment{theoreme}{leftlined={5pt}{red!40!darkgray}}
\theoremstyle{theoreme}
\newtheorem{theoreme}[proposition]{Théorème}
%\newtcbtheorem[number within=chapter]{proposition}{Proposition}





%-------------------
%---- Lemmes ----
%-------------------
\newtheoremstyle{lemme}% name
{3pt}% Space above
{1cm}% Space below
{\itshape}% Body font
{}% Indent amount
{\bfseries}% Theorem head font
{\vspace{1mm}\newline}% Punctuation after theorem head
{1cm}%{0pt\Needspace{3\baselineskip}}% Space after theorem head
{}% Theorem head spec (can be left empty, meaning ‘normal’ )

\tcolorboxenvironment{lemme}{leftlined={5pt}{gray}}
\theoremstyle{lemme}
\newtheorem{lemme}[proposition]{Lemme}
%\newtcbtheorem[number within=chapter]{proposition}{Proposition}







%-----------------------
%---- Démonstration ----
%-----------------------
\newtheoremstyle{demonstration}% name
{3pt}% Space above  % -> NE MARCHE PAS - CERTAINEMENT A CAUSE DE TCB QUI DOIT "OVERRULER" LES DIMENSIONS
{3pt}% Space below  % -> NE MARCHE PAS - CERTAINEMENT A CAUSE DE TCB QUI DOIT "OVERRULER" LES DIMENSIONS
{}% Body font \itshape
{-10mm}% Indent amount
{\itshape}% Theorem head font
{\vspace{2.5mm}\newline}% Punctuation after theorem head
{1mm}%{0pt\Needspace{3\baselineskip}}% Space after theorem head % -> NE MARCHE PAS - CERTAINEMENT A CAUSE DE TCB QUI DOIT "OVERRULER" LES DIMENSIONS
{}% Theorem head spec (can be left empty, meaning ‘normal’ )

\theoremstyle{demonstration}
\tcolorboxenvironment{demonstration}{indented={8mm}}
\newtheorem*{demonstration}{Démonstration :}







%------------------
%---- Exercice ----
%------------------
\tikzset{
     shadowed/.style={preaction={transform canvas={shift={(-1pt,-1pt)}},draw=gray,very thick}},
  }
\newcounter{exercices}
% colexam est définit dans le préambule du document maître
\newtcolorbox[use counter=exercices]{exstheorie}{
empty,
coltext=gray!10!black,
title={Exercices},
attach boxed title to top left,
boxed title style={toprule=2pt,top=2pt,left=-0.6cm,
overlay={\path[left color=white,right color=colexam,rounded corners=5pt]([yshift=0pt]frame.north east)+(0cm,-23pt) rectangle +(-5cm,0pt);}},
overlay unbroken={\draw[colexam,line width=1pt,shadowed]
([xshift=-10mm,yshift=-1pt]title.south-|frame.west)
--([xshift=-10mm]frame.south west)--(frame.south); },
coltitle=white,fonttitle=\Large\bfseries\itshape,
before=\par\medskip\noindent,parbox=false,boxsep=0pt,left=-5mm,right=3mm,top=2mm,bottom=5mm,
%breakable,pad at break*=0mm,vfill before first,
%overlay first={\draw[colexam,line width=1pt]
%([yshift=-1pt]title.north east)--([xshift=-0.5pt,yshift=-1pt]title.north-|frame.east)
%--([xshift=-0.5pt]frame.south east); }%,
%overlay middle={\draw[colexam,line width=1pt] ([xshift=-0.5pt]frame.north east)
%--([xshift=-0.5pt]frame.south east); }%,
%overlay last={\draw[colexam,line width=1pt] ([xshift=-0.5pt]frame.north east)
%--([xshift=-0.5pt]frame.south east)--(frame.south west);}%,%
}


%\newcounter{exercices}
%\colorlet{colexam}{gray!75!black}
%\newtcolorbox[use counter=exercices]{extheorie}{%
%empty,title={Exercices},attach boxed title to top left,
%boxed title style={empty,size=minimal,toprule=2pt,top=6pt,overlay={\draw[colexam,line width=2.5pt]([yshift=-1pt]frame.north west)--([yshift=-1pt]frame.north east);}},
%coltitle=black,fonttitle=\Large\bfseries\itshape,
%before=\par\medskip\noindent,parbox=false,boxsep=0pt,left=10pt,right=3mm,top=10pt,
%breakable,pad at break*=0mm,vfill before first,
%overlay unbroken={\draw[colexam,line width=1pt]
%([yshift=-1pt]title.north east)--([xshift=-0.5pt,yshift=-1pt]title.north-|frame.east)
%--([xshift=-0.5pt]frame.south east)--(frame.south west); },
%%overlay first={\draw[colexam,line width=1pt]
%%([yshift=-1pt]title.north east)--([xshift=-0.5pt,yshift=-1pt]title.north-|frame.east)
%%--([xshift=-0.5pt]frame.south east); }%,
%%overlay middle={\draw[colexam,line width=1pt] ([xshift=-0.5pt]frame.north east)
%%--([xshift=-0.5pt]frame.south east); }%,
%%overlay last={\draw[colexam,line width=1pt] ([xshift=-0.5pt]frame.north east)
%%--([xshift=-0.5pt]frame.south east)--(frame.south west);}%,%
%}


%%%%%%%%%%%%%%%%%%%%%%%%%%%%%%%%
%D'autres possibilités pour les exercices. A travailler.
%%%%%%%%%%%%%%%%%%%%%%%%%%%%%%%%
%
%\begin{tcolorbox}[enhanced,colframe=gray!75!white,interior style={left color=gray!30,right color=white},leftrule=3mm,title=Exercices]
%This is the upper part.
%
%This is the lower part.
%\end{tcolorbox}
%
%%%%%%%%%%%%%%%%%%%%%%%%%%%%%%%%
%
%\begin{tcolorbox}[enhanced,sharp corners=south,
%colframe=red!75!black,interior style={top color=white,bottom color=red!50}]
%This is an example for 'spread downwards'.
%\end{tcolorbox}
%
%%%%%%%%%%%%%%%%%%%%%%%%%%%%%%%%
%
%\newtheoremstyle{exercice}% name
%{3pt}% Space above  % -> NE MARCHE PAS - CERTAINEMENT A CAUSE DE TCB QUI DOIT "OVERRULER" LES DIMENSIONS
%{3pt}% Space below  % -> NE MARCHE PAS - CERTAINEMENT A CAUSE DE TCB QUI DOIT "OVERRULER" LES DIMENSIONS
%{\itshape}% Body font \itshape
%{-2mm}% Indent amount
%{\bfseries\itshape}% Theorem head font
%{\newline}% Punctuation after theorem head
%{1mm}%{0pt\Needspace{3\baselineskip}}% Space after theorem head % -> NE MARCHE PAS - CERTAINEMENT A CAUSE DE TCB QUI DOIT "OVERRULER" LES DIMENSIONS
%{}% Theorem head spec (can be left empty, meaning ‘normal’ )
%
%\theoremstyle{exercice}
%\tcolorboxenvironment{exercice}{exstyle}
%\newtheorem*{exercice}{Exercice :}
%
%
%%-------------------
%%---- ExerciceS ----
%%-------------------
%\newtheoremstyle{exercices}% name
%{3pt}% Space above  % -> NE MARCHE PAS - CERTAINEMENT A CAUSE DE TCB QUI DOIT "OVERRULER" LES DIMENSIONS
%{3pt}% Space below  % -> NE MARCHE PAS - CERTAINEMENT A CAUSE DE TCB QUI DOIT "OVERRULER" LES DIMENSIONS
%{\itshape}% Body font \itshape
%{-2mm}% Indent amount
%{\bfseries\itshape}% Theorem head font
%{\newline}% Punctuation after theorem head
%{1mm}%{0pt\Needspace{3\baselineskip}}% Space after theorem head % -> NE MARCHE PAS - CERTAINEMENT A CAUSE DE TCB QUI DOIT "OVERRULER" LES DIMENSIONS
%{}% Theorem head spec (can be left empty, meaning ‘normal’ )
%
%\theoremstyle{exercices}
%\tcolorboxenvironment{exercices}{indented={1mm}}
%\newtheorem*{exercices}{Exercices :}








%%%%%%%%%%%%% ANCIENNE VERSION %%%%%%%%%%%%%
%\newenvironment{prop}[1]
%	{
%		\refstepcounter{compteurthm}         
%         \begin{center}%
%            \begin{tikzpicture}%
%                \node[rectangle, draw=black, rounded corners=0pt, inner xsep=5pt, inner ysep=6pt, outer ysep=10pt]\bgroup                     
%                		\begin{minipage}{#1\linewidth}
%                			{\bf Proposition \thecompteurthm\hspace{2pt}:}
%                			
%                			\vspace{2mm}\hspace{7.5mm}
%                			\begin{minipage}{0.93\linewidth}\it
%
%    }
%	{                		
%                			\end{minipage}
%                		\end{minipage}\egroup;%
%            \end{tikzpicture}%
%         \end{center}%
%	}




%------------------------
%---- Démonstrations ----
%------------------------
\newenvironment{demo}[1]
	{
        \begin{center}%
			\begin{minipage}{#1\linewidth}
				\textit{Démonstration :}
				
				\vspace{1mm}
    }
	{
			
			%\vspace{-4mm}
			\hfill \textit{CQFD}\hspace{1mm} $\blacksquare$
        	\end{minipage}
        \end{center}
	}




%-----------------------------------------------
%---- RemarqueS : en encadré sur fond blanc ----
%-----------------------------------------------
\newcounter{compteuremarque}
\setcounter{compteuremarque}{1}
\newenvironment{remarques}
	{
%         \begin{center}%
%            \begin{tikzpicture}%
%                \node[rectangle, draw=black, rounded corners=5pt, inner xsep=5pt, inner ysep=6pt, outer ysep=10pt]\bgroup                     
%                		\begin{minipage}{0.97\linewidth}
                			\textbf{Remarques \thecompteuremarque.\;}
                			\vspace*{-1mm}
                			\begin{enumerate}[label=\roman*), itemsep=0mm]%\setlength{\itemsep}{1mm}
    }
	{                		
                			\end{enumerate}\stepcounter{compteuremarque}
%                		\end{minipage}%\egroup;%
%            \end{tikzpicture}%
%         \end{center}%
	}




%----------------------------------------------
%---- Remarque : en encadré sur fond blanc ----
%----------------------------------------------
\newenvironment{remarque}
	{
%         \begin{center}%
%            \begin{tikzpicture}%
%                \node[rectangle, draw=black, rounded corners=5pt, inner xsep=5pt, inner ysep=6pt, outer ysep=10pt]\bgroup                     
%                		\begin{minipage}{0.97\linewidth}
                			{\bf Remarque \thecompteuremarque.}
                			
                			\vspace{2mm}\hspace{7.5mm}
                			\begin{minipage}{0.93\linewidth}

    }
	{                		
                			\end{minipage}\stepcounter{compteuremarque}
%                		\end{minipage}\egroup;%
%            \end{tikzpicture}%
%         \end{center}%
	}

















%-------------------
%---- Exemple : ----
%-------------------
\newenvironment{exemple}
	{
	\noindent\textit{Exemple :}\hspace{2mm}\begin{minipage}[t]{0.9\textwidth}
	}  
	{
	\end{minipage}
	}




%%--------------------------
%%---- ExempleS : Liste ----
%%--------------------------
%\newenvironment{exemples}
%	{
%	\noindent\textit{Exemples :}\begin{minipage}[t][][b]{0.9\textwidth}\begin{enumerate}[label=\arabic*)]
%	}  
%	{
%	\end{enumerate}
%	\end{minipage}
%	}
%--------------------------
%---- ExempleS : Liste ----
%--------------------------
\newenvironment{exemples}[1]
	{
	\noindent\textit{Exemples :}\bgroup
	\vspace*{-1mm}\begin{enumerate}[label=\arabic*),itemsep=#1]
	}  
	{
	\end{enumerate}\egroup
	}




%%--------------------------------
%%---- Exercice : en italique ----
%%--------------------------------
%\newenvironment{exercice}
%	{
%	\noindent\textit{\textbf{Exercice:}} \begin{minipage}[t]{0.9\textwidth}\itshape
%	}  
%	{
%	\end{minipage}
%	}




%---------------------------------------------------------------------------------
%---- Attention : en encadré sur fond rouge avec symbôle "attention" à gauche ----
%---------------------------------------------------------------------------------
\definecolor{colimportant}{RGB}{247 , 189 , 164}
\definecolor{contourimportant}{RGB}{237 , 125 , 49}
\newenvironment{important}[2]
	{
         \begin{center}%
            \begin{tikzpicture}%
                \node[rectangle, draw=black!80, top color=white, bottom color=white, rounded corners=5pt, inner xsep=5pt, inner ysep=6pt, outer ysep=10pt]\bgroup                     
                		\begin{minipage}{0.08\linewidth}\centerline{\includegraphics[scale=#2]{#1}}\end{minipage}
                		\begin{minipage}{0.89\linewidth}
    }
	{                		
                		\end{minipage}\egroup;
            \end{tikzpicture}%
         \end{center}%
	}




%-------------------------------------------------------------------------------
%---- Eclairage : en encadré sur fond jaune avec symbôle "ampoule" à gauche ----
%-------------------------------------------------------------------------------
\definecolor{coleclairage}{RGB}{255 , 221 , 156}
\definecolor{contoureclairage}{RGB}{255 , 192 , 0}
\newenvironment{eclairage}
	{\vspace{-7mm}
         \begin{center}%
            \begin{tikzpicture}%
                \node[rectangle, draw=contoureclairage, top color=coleclairage!50, bottom color=coleclairage!140, rounded corners=5pt, inner xsep=5pt, inner ysep=6pt, outer ysep=10pt]\bgroup                     
                		\begin{minipage}{0.08\linewidth}\centerline{\includegraphics[scale=1]{Symbole_eclairage.png}}\end{minipage}
                		\begin{minipage}{0.89\linewidth}\comicneue\itshape\footnotesize
    }
	{                		
                		\end{minipage}\egroup;%
            \end{tikzpicture}%
         \end{center}%
     \vspace{-1mm}
	}




%-----------------------------------------------------------------
%---- Choses semi-importantes (formules ou autres) en encadré ----
%-----------------------------------------------------------------
\newenvironment{encadre}[1]
	{
         \begin{center}%
            \begin{tikzpicture}%
                \node[rectangle, draw=black, rounded corners=1pt, inner xsep=5pt, inner ysep=6pt, outer ysep=10pt]\bgroup                     
                		\begin{minipage}{#1\linewidth}
							\begin{center}
    }
	{
							\end{center}
                		\end{minipage}\egroup;%
            \end{tikzpicture}%
         \end{center}%
	}



%-----------------------------------------------------
%---- Affichage PYTHON : en encadré sur fond gris ----
%-----------------------------------------------------
\newenvironment{affichage}[1]
	{
         \begin{center}%
            \begin{tikzpicture}%
                \node[rectangle, draw=black, fill=lightgray, rounded corners=5pt, inner xsep=5pt, inner ysep=6pt, outer ysep=10pt]\bgroup                     
                		\begin{minipage}{#1\linewidth}

    }
	{                		
                		\end{minipage}\egroup;%
            \end{tikzpicture}%
         \end{center}%
	}



%%%%%%%%%%%%%%%%%%%%%%%%%%%%%%%%%%%%%%%%%%%%%%%%%%%%%%%%%%%%%%%%%%%%%%%%%%%%%%%%%%%%%%%%%%%%%%%%%%%%%%%%%%%%%%%%%%%%%%%%%%%%%%%%%%%%%%%%%
%%%%%%%%%%%%%%%%%%%%%%%%%%%%   DEFINITIONS ALTERNATIVES (QUI FONCTIONNENT) POUR LES ENVIRONNEMENTS ENCADRES   %%%%%%%%%%%%%%%%%%%%%%%%%%%
%%%%%%%%%%%%%%%%%%%%%%%%%%%%%%%%%%%%%%%%%%%%%%%%%%%%%%%%%%%%%%%%%%%%%%%%%%%%%%%%%%%%%%%%%%%%%%%%%%%%%%%%%%%%%%%%%%%%%%%%%%%%%%%%%%%%%%%%%
%     
%%%%%%%%%%%%%%%%%%%%%%%%%%%%  Définition avec l'implémentation du texte à afficher dans une boîte  %%%%%%%%%%%%%%%%%%%%%%%%%%%
%\newenvironment{important}
%	{%
%	\setbox0=\hbox\bgroup % On génére une boîte qui s'appelle "box0" et on met le contenu à afficher dans cette boîte.       
%		\begin{minipage}{0.88\textwidth}	
%	}         
%    {
%		\end{minipage}    
%    	\egroup % On ferme la boîte
%    	\begin{center}% On commende à construire l'environnement
%            \begin{tikzpicture}%
%                \node[rectangle, draw=contourimportant, top color=colimportant!50, bottom color=colimportant!140, rounded corners=5pt, inner xsep=5pt, inner ysep=6pt, outer ysep=10pt]{                      
%                		\begin{minipage}{0.08\textwidth}\centerline{\includegraphics[scale=1]{Symbole_attention.png}}\end{minipage}
%                		\begin{minipage}{0.89\textwidth} \box0	\end{minipage} }; % On introduit la boîte dans laquelle il y a le texte à afficher.
%            \end{tikzpicture}%
%         \end{center}%
%	}
%
%
%%%%%%%%%%%%%%%%%%%%%%%%%%%%  Définition dans une commande paramétrée  %%%%%%%%%%%%%%%%%%%%%%%%%%%
%\newcommand{\important}[1]{%
%         \begin{center}%
%            \begin{tikzpicture}%
%                \node[rectangle, draw=contourimportant, top color=colimportant!50, bottom color=colimportant!140, rounded corners=5pt, inner xsep=5pt, inner ysep=6pt, outer ysep=10pt]{                      
%                		\begin{minipage}{0.08\linewidth}\centerline{\includegraphics[scale=1]{Symbole_attention.png}}\end{minipage}
%                		\begin{minipage}{0.89\linewidth}#1\end{minipage}};%
%            \end{tikzpicture}%
%         \end{center}%






%%%%%%%%%%%%%%%%%%%%%%%%%%%%%%%%%%%%%%%%%%%%%%%%%%%%%%%%%%%%%%%%%%%%%%%%%%%%%%%%%
%%%%%%%%%%%%%%%%%%%%%%%%%%%   COMMANDES PARAMETREES   %%%%%%%%%%%%%%%%%%%%%%%%%%%
%%%%%%%%%%%%%%%%%%%%%%%%%%%%%%%%%%%%%%%%%%%%%%%%%%%%%%%%%%%%%%%%%%%%%%%%%%%%%%%%%
%\newcommand{\fig}[1]{\begin{center} {\textbf{\textit{ Figure \thechapter.\thenfigure : }}}{\footnotesize{\it #1}\stepcounter{nfigure}}\end{center}} %la commande "the(compteur)" permet d'afficher ou en est le comteur.
%\newenvironment{remarque}[1]{{\bf Remarque \thechapter.\thecompteurthm\;: }\\ \hspace*{1.1cm}\begin{minipage}{0.9\textwidth}{\vspace*{2mm} #1}\end{minipage}\stepcounter{compteurthm}}

%%---- Choses semi-importantes (formules ou autres) en encadré ----
%\newenvironment{encadre}[2]{\begin{center}\fbox{\begin{minipage}{#1\textwidth}{\begin{center}#2\end{center}}\end{minipage}}\end{center}}

\newcounter{nfigure}%[chapter] % Ici on définit un nouveau compteur associé au compteur "chapter", i.e. quand "chapter" augmente de niveau, "nfigure" est remis à zéro.
%\newcommand{\fig}[1]{\begin{center} {\textbf{\textit{ Figure \thechapter.\thenfigure : }}}{\footnotesize{\it #1}\stepcounter{nfigure}}\end{center}} %la commande "the(compteur)" permet d'afficher ou en est le comteur.
\setcounter{nfigure}{1}
\newcommand{\fig}{Figure \thechapter.\thenfigure \stepcounter{nfigure}}






\newenvironment{changemargin}[2]{\begin{list}{}{%
\setlength{\topsep}{0pt}%
\setlength{\leftmargin}{0pt}%
\setlength{\rightmargin}{0pt}%
\setlength{\listparindent}{\parindent}%
\setlength{\itemindent}{\parindent}%
\setlength{\parsep}{0pt plus 1pt}%
\addtolength{\leftmargin}{#1}%
\addtolength{\rightmargin}{#2}%
}\item }{\end{list}}






